\subsection{\texorpdfstring{\(\mathfrak{sl}(2,k)\) -模的本原元}{\textbackslash mathfrak\{sl\}(2,k) -模的本原元}}

设 E 是 \(\mathfrak{sl}(2,k)\) -模。若 \(A \in \mathfrak{sl}(2,k)\) 且 \(x \in E\) ，我们常把 \(A_{E}x\) 写作 Ax。设 \(\lambda \in k\) 。若 Hx = \(\lambda x\) ，则按语言滥用，我们称 x 是 E 中权为 \(\lambda\) 的元素，或称 \(\lambda\) 是 x 的权。若 E 有限维，则 \(H_{E}\) 半单，故权为 \(\lambda\) 的元素之集是 E 相对于 \(H_{E}\) 与 \(\lambda\) 的准素子空间（参见第七章 §1 第 1 节）。

引理 3. 若 \(x\) 是权为 \(\lambda\) 的元素，则 \(X_{+}x\) 是权为 \(\lambda + 2\) 的元素，而 \(X_{-}x\) 是权为 \(\lambda - 2\) 的元素。

事实上，\(HX_{+}x = [H,X_{+}]x + X_{+}Hx = 2X_{+}x + X_{+}\lambda x = (\lambda +2)X_{+}x\) ，同理 \(HX_{-}x = (\lambda -2)X_{-}x\) （另见第七章 §1 第 3 节 命题 10 (ii)）。

定义 1. 设 E 是 \(\mathfrak{sl}(2,k)\) -模。E 的元素称为本原元，如果它是 \(H_{E}\) 的非零特征向量且属于 \(X_{+E}\) 的核。

E 的非零元素 e 是本原元，当且仅当 ke 在 \(kH + kX_{+}\) 的作用下稳定；这例如可由引理 3 得出。

例 元素 \(X_{+}\) 对 \(\mathfrak{sl}(2,k)\) 的伴随表示是权 2 的本原元。元素 \((1,0)\) （\(k^{2}\) 中的）对 \(\mathfrak{sl}(2,k)\) 在 \(k^{2}\) 上的恒同表示是权 1 的本原元。

引理 4. 设 \(\mathbf{E}\) 是非零有限维 \(\mathfrak{sl}(2,k)\) -模。则 \(\mathbf{E}\) 有本原元。

由于 \(X_{+}\) 是 \(\mathfrak{sl}(2,k)\) 的幂零元素，\(X_{+E}\) 幂零。假设 \(X_{+E}^{m-1} \neq 0\) 且 \(X_{+E}^{m} = 0\) 。由引理 2，

\[
m (H _ {\mathrm{E}} - m + 1) X _ {+ \mathrm{E}} ^ {m - 1} = [ X _ {- \mathrm{E}}, X _ {+ \mathrm{E}} ^ {m} ] = 0,
\]

因而 \(X_{+}^{m - 1}(\mathrm{E}) - \{0\}\) 中的元素都是本原的。

命题 1. 设 E 是 \(\mathfrak{sl}(2,k)\) -模，\(e\) 是 E 的权为 \(\lambda\) 的本原元。置 \(e_n = \frac{(-1)^n}{n} X_-^n e\) （对 \(n \geq 0\) ），并置 \(e_{-1} = 0\) 。则

\[
\left\{ \begin{array}{l l} H e _ {n} & = (\lambda - 2 n) e _ {n} \\ X _ {-} e _ {n} & = - (n + 1) e _ {n + 1} \\ X _ {+} e _ {n} & = (\lambda - n + 1) e _ {n - 1}. \end{array} \right.\tag{2}
\]

第一个公式由引理 3 得出，第二个由 \(e_n\) 的定义得出。我们对 \(n\) 作归纳来证明第三个。它在 \(n = 0\) 时成立，因为 \(e_{-1} = 0\) 。若 \(n > 0\) ，

\[
\begin{array}{r l} & n X _ {+} e _ {n} = - X _ {+} X _ {-} e _ {n - 1} = - [ X _ {+}, X _ {-} ] e _ {n - 1} - X _ {-} X _ {+} e _ {n - 1} \\ & \qquad = H e _ {n - 1} - X _ {-} (\lambda - n + 2) e _ {n - 2} \\ & \qquad = (\lambda - 2 n + 2 + (n - 1) (\lambda - n + 2)) e _ {n - 1} \\ & \qquad = n (\lambda - n + 1) e _ {n - 1}. \end{array}
\]

推论. \(\mathbf{E}\) 的由 \(e\) 生成的子模是由诸 \(e_n\) 生成的向量子空间。

这由公式 (2) 得出。

整数 \(n \geq 0\) 中使 \(e_{n} \neq 0\) 者构成 N 中一个区间，而相应的元素 \(e_{n}\) 在 k 上构成由 e 生成的子模的基（事实上，它们线性无关，因为它们是互异权的非零元素）。称此基为与本原元 e 相伴的基。

命题 2. 若由本原元 e 生成的 E 的子模 V 是有限维的，则：

\begin{enumerate}
\def\labelenumi{(\roman{enumi})}
\item
  e 的权 \(\lambda\) 是整数且等于 \(\dim V - 1\) ；
\item
  \((e_0, e_1, \ldots, e_\lambda)\) 是 V 的基，且 \(e_n = 0\) 对 \(n > \lambda\) 成立；
\item
  \(H_{\mathrm{V}}\) 的特征值是 \(\lambda, \lambda - 2, \lambda - 4, \ldots, -\lambda\) ；它们的重数均为 1；
\item
  V 的每个本原元都与 e 成比例；
\item
  模 \(\mathrm{V}\) 的交换子可约化为纯量；特别地，\(\mathrm{V}\) 绝对单。
\end{enumerate}

设 \(m\) 是使 \(e_m \neq 0\) 的最大整数。则 \(0 = X_{+}e_{m+1} = (\lambda - m)e_m\) ，故 \(\lambda = m\) ；由于 \((e_0, e_1, \ldots, e_m)\) 是 \(V\) 的基，这就证明了 (i) 与 (ii)。断言 (iii) 由等式 \(He_n = (\lambda - 2n)e_n\) 得出。我们有 \(X_{+}e_n \neq 0\) 对 \(1 \leq n \leq m\) 成立，故得 (iv)。设 \(c\) 是模 \(V\) 的交换子的元素。则 \(Hc(e) = cH(e) = \lambda c(e)\) ，故存在 \(\mu \in k\) 使得 \(c(e) = \mu e\) ；于是

\[
c X _ {-} ^ {q} e = X _ {-} ^ {q} c e = \mu X _ {-} ^ {q} e
\]

对一切 \(q \geq 0\) 成立，故 \(c = \mu.1\) ，这就证明了 (v)。

推论. 设 E 是有限维 \(\mathfrak{sl}(2,k)\) -模。

\begin{enumerate}
\def\labelenumi{(\roman{enumi})}
\item
  自同态 \(H_{E}\) 可对角化，其特征值是有理整数。
\item
  对任何 \(p \in Z\) ，设 \(E_{p}\) 是 \(H_{E}\) 的对应于特征值 p 的特征子空间。设 i 是整数 \(\geq 0\) 。映射 \(X_{-E}^{i}|E_{p}:E_{p} \to E_{p-2i}\) 在 \(i \leq p\) 时单，在 i = p 时双，在 \(i \geq p\) 时满。映射 \(X_{+E}^{i}|E_{-p}:E_{-p} \to E_{-p+2i}\) 在 \(i \leq p\) 时单，在 i = p 时双，在 \(i \geq p\) 时满。
\item
  \(\mathrm{E}\) 的长度等于 \(\dim \operatorname{Ker} X_{+\mathrm{E}}\) ，也等于 \(\dim \operatorname{Ker} X_{-\mathrm{E}}\) 。
\item
  设 \(E'\) （相应地 \(E''\) ）是对 p 为偶（相应地奇）的诸 \(E_{p}\) 之和。则 \(E'\) （相应地 \(E''\) ）是 E 的奇（相应地偶）维单子模之和；且 \(E = E' \oplus E''\) 。\(E'\) 的长度是 \(\dim E_{0}\) ，\(E''\) 的是 \(\dim E_{1}\) 。
\end{enumerate}

\[
\text {(v)} \operatorname{Ker} X _ {+ \mathrm{E}} \cap \operatorname{Im} X _ {+ \mathrm{E}} \subset \sum_ {p > 0} \mathrm{E} _ {p} \text {and} \operatorname{Ker} X _ {- \mathrm{E}} \cap \operatorname{Im} X _ {- \mathrm{E}} \subset \sum_ {p <   0} \mathrm{E} _ {p}.
\]

若 E 单，则 E 由一个本原元生成（引理 4），只需应用命题 1 与 2。一般情形由此推出，因为每个有限维 \(\mathfrak{sl}(2,k)\) -模都半单。

